\documentclass[letterpaper]{article}
\usepackage{aaai2027}
\usepackage{amsmath,amssymb}
\usepackage{booktabs}
\usepackage{graphicx}
\usepackage{multirow}
\usepackage{xcolor}
\usepackage{url}

\newcommand{\method}{RPC-CLIP}
\newcommand{\todo}[1]{\textcolor{red}{[#1]}}

\title{Reliability-Guided Prompt-Prototype Calibration for\\Weakly Supervised Semantic Segmentation}
\author{Anonymous AAAI 2027 Submission}

\begin{document}
\maketitle

\begin{abstract}
Weakly supervised semantic segmentation (WSSS) seeks to learn dense masks from
image-level labels, but modern CLIP-based pipelines still suffer from noisy
patch-text activations, incomplete object coverage, and unstable background
assignment. We propose \method, a single-GPU friendly framework that freezes
CLIP and trains only a lightweight patch decoder. The key idea is to calibrate
CLIP pseudo labels with a momentum prototype bank and an entropy-based
reliability estimator. Given image-level labels, \method first derives
class-constrained patch scores from prompt-averaged CLIP text embeddings. It then
updates foreground and background prototypes from high-confidence patches, fuses
text, prototype, and decoder distributions according to reliability, and applies
a boundary-aware smoothness regularizer. This design keeps training fast while
progressively correcting under-activated object regions. Experiments on PASCAL
VOC 2012 and optional MS COCO will show that \method improves over text-only
CLIP pseudo labels and remains competitive with recent CLIP-based WSSS methods.
\todo{Fill final mIoU, test server numbers, runtime, and significance after
real experiments.}
\end{abstract}

\section{Introduction}
Weakly supervised semantic segmentation reduces annotation cost by replacing
pixel-level masks with cheaper supervision such as image-level labels. Classic
CAM-based methods~\cite{zhou2016cam} localize discriminative object parts, but
they often miss complete object extents. Recent CLIP-based
methods~\cite{radford2021clip,lin2023clipes,zhang2024weclip} offer stronger
open-vocabulary semantic priors, yet patch-text similarity maps remain noisy:
confident patches may cover only the most recognizable parts, while background
pixels can be mistakenly attracted by foreground prompts.

This paper asks whether a frozen CLIP model can be made more reliable without
introducing a heavy segmentation backbone. Our answer is \method, a
prompt-prototype calibration framework designed for fast single-GPU training.
Instead of fine-tuning CLIP, we keep the visual-language encoder fixed and train
a small patch decoder on pseudo labels. The pseudo labels are not taken directly
from CLIP. They are calibrated by a momentum prototype bank that accumulates
class evidence from confident patches and by a reliability score that suppresses
uncertain regions.

The proposed framework follows a practical principle: a pseudo label should be
trusted only when language evidence, visual prototypes, and the current decoder
are sufficiently consistent. This reliability-guided fusion improves training
stability and gives the decoder a cleaner curriculum. To reduce oversmoothing
across object boundaries, we further add an image-gradient weighted smoothness
term.

Our contributions are:
\begin{itemize}
  \item We introduce a lightweight frozen-CLIP WSSS framework that trains on a
  single GPU and uses only image-level labels.
  \item We propose prompt-prototype calibration, where class-constrained CLIP
  text scores are corrected by a momentum prototype bank.
  \item We design a reliability-guided pseudo-label fusion objective combining
  text, prototype, and decoder distributions.
  \item We provide a complete experimental protocol with comparison, ablation,
  and visualization templates for PASCAL VOC 2012.
\end{itemize}

\section{Related Work}
\textbf{Weakly supervised semantic segmentation.}
CAM-based WSSS methods use image classifiers to localize discriminative regions.
SEAM~\cite{wang2020seam} improves consistency under image transformations,
while reliability-aware methods~\cite{zhang2021rrm,zhang2021a2gnn} show that
mining trustworthy regions is crucial for end-to-end WSSS. \method follows this
reliability-oriented line but replaces a supervised classification backbone with
a frozen CLIP encoder.

\textbf{Vision-language models for segmentation.}
CLIP~\cite{radford2021clip} transfers language supervision to visual recognition
and has become a strong prior for weakly supervised segmentation. CLIP-ES and
WeCLIP~\cite{lin2023clipes,zhang2024weclip} demonstrate that CLIP patch-text
similarities can produce useful pseudo masks. However, raw language similarity
is not calibrated for dense prediction. Our method uses text features as
initialization and continuously refines them with image-level visual prototypes.

\textbf{Prototype learning.}
Prototype-based segmentation summarizes class appearance with representative
features. SFC~\cite{xu2024sfc} shows that feature calibration can improve weakly
supervised segmentation. In weak supervision, prototype updates must be
conservative because incorrect pseudo labels can quickly contaminate the class
representation. \method addresses this issue with confidence-gated momentum
updates and entropy-based reliability weighting.

\section{Method}
\subsection{Overview}
Given an image $I$ and image-level label vector $y \in \{0,1\}^{C}$, a frozen
CLIP ViT extracts patch features $X=\{x_i\}_{i=1}^{N}$, where
$x_i \in \mathbb{R}^{d}$. A lightweight decoder $g_\theta$ predicts dense logits
$Z=g_\theta(X) \in \mathbb{R}^{(C+1)\times H \times W}$ including background.
Only $\theta$ is optimized.

\subsection{Prompt-Constrained CLIP Seeds}
For each class $c$, we average CLIP text embeddings over prompt templates:
\begin{equation}
t_c = \mathrm{norm}\left(\frac{1}{K}\sum_{k=1}^{K}
\mathrm{CLIP}_{text}(p_k(c))\right).
\end{equation}
Patch-text logits are computed as cosine similarities:
\begin{equation}
s^{text}_{i,c} = x_i^\top t_c / \tau_t.
\end{equation}
Classes absent from $y$ are masked out. Background is initialized as the inverse
of the strongest foreground response, encouraging low-foreground patches to be
used as background candidates.

\subsection{Momentum Prototype Calibration}
We maintain prototypes $P=\{p_0,\ldots,p_C\}$ for background and foreground
classes. Foreground prototypes are initialized by text embeddings, while the
background prototype is learned from confident low-foreground regions. Given
pseudo target $\hat{m}$ and confidence $q$, class $c$ is updated only when enough
patches satisfy $\hat{m}_i=c$ and $q_i>\gamma$:
\begin{equation}
p_c \leftarrow \mathrm{norm}\left(\mu p_c + (1-\mu)
\frac{1}{|\Omega_c|}\sum_{i\in\Omega_c} x_i\right).
\end{equation}
Prototype logits are $s^{proto}_{i,c}=x_i^\top p_c/\tau_p$.

\subsection{Reliability-Guided Fusion}
The final pseudo distribution combines text, prototype, and decoder
probabilities:
\begin{equation}
\pi_i = \lambda_t \sigma(s^{text}_i) + \lambda_p \sigma(s^{proto}_i)
+ \lambda_d \sigma(z_i).
\end{equation}
Reliability is estimated by normalized entropy and confidence:
\begin{equation}
r_i = \max_c \pi_{i,c}\left(1 -
\frac{-\sum_c \pi_{i,c}\log\pi_{i,c}}{\log(C+1)}\right).
\end{equation}
Pixels whose confidence is lower than $\eta$ are ignored. The training objective
is:
\begin{equation}
\mathcal{L}= \mathcal{L}_{ce} + \alpha\mathcal{L}_{cls}
+ \beta\mathcal{L}_{kl} + \delta\mathcal{L}_{bd}.
\end{equation}
Here $\mathcal{L}_{ce}$ is reliability-weighted cross entropy,
$\mathcal{L}_{cls}$ is image-level multi-label classification loss from pooled
decoder logits, $\mathcal{L}_{kl}$ aligns decoder probabilities to $\pi$, and
$\mathcal{L}_{bd}$ is an image-gradient weighted smoothness term.

\section{Experiments}
\subsection{Datasets and Metrics}
We evaluate on PASCAL VOC 2012~\cite{voc2010} using image-level labels for
training and report mean intersection-over-union (mIoU) on validation and test
sets. Optional MS COCO experiments can be added after the VOC pipeline is
stable.

\subsection{Implementation Details}
Unless specified otherwise, \method uses CLIP ViT-B/16, four prompt templates,
input resolution $224\times224$, AdamW, batch size 16, and eight fast
pre-experiment epochs. The CLIP encoder is frozen. Only the decoder is trained.
All experiments are run on a single GPU. \todo{Replace with exact GPU model and
training time.}

\subsection{Comparison with State of the Art}
\begin{table}[t]
\centering
\small
\begin{tabular}{lccc}
\toprule
Method & Backbone & Val mIoU & Test mIoU \\
\midrule
CAM & VGG16 & \todo{TBD} & \todo{TBD} \\
SEAM & ResNet38 & \todo{TBD} & \todo{TBD} \\
RRM & ResNet38 & \todo{TBD} & \todo{TBD} \\
CLIP-ES & CLIP ViT-B/16 & \todo{TBD} & \todo{TBD} \\
SFC & CLIP ViT-B/16 & \todo{TBD} & \todo{TBD} \\
WeCLIP & CLIP ViT-B/16 & \todo{TBD} & \todo{TBD} \\
\method & CLIP ViT-B/16 & \todo{TBD} & \todo{TBD} \\
\bottomrule
\end{tabular}
\caption{Comparison on PASCAL VOC 2012. Fill all numbers after reproducing or
quoting under the same protocol.}
\end{table}

\subsection{Ablation Study}
\begin{table}[t]
\centering
\small
\begin{tabular}{lcc}
\toprule
Variant & Val mIoU & $\Delta$ \\
\midrule
Text-only pseudo labels & \todo{TBD} & -- \\
+ Prototype calibration & \todo{TBD} & \todo{TBD} \\
+ Reliability weighting & \todo{TBD} & \todo{TBD} \\
+ Boundary smoothness (\method) & \todo{TBD} & \todo{TBD} \\
\bottomrule
\end{tabular}
\caption{Ablation of RPC components.}
\end{table}

\subsection{Qualitative Results}
Figure~\ref{fig:qualitative} will show input images, pseudo labels, predictions,
and ground-truth masks. We will include both success and failure cases to reveal
whether reliability calibration improves object coverage.
\begin{figure}[t]
\centering
\fbox{\rule{0pt}{1.2in}\rule{0.9\linewidth}{0pt}}
\caption{Qualitative examples. Replace this placeholder with panels generated
by \texttt{visualize.py}.}
\label{fig:qualitative}
\end{figure}

\section{Conclusion}
We presented \method, a fast frozen-CLIP framework for weakly supervised
semantic segmentation. By combining prompt-constrained language evidence,
momentum prototype calibration, and reliability-guided pseudo-label fusion,
\method aims to improve pseudo-mask quality while keeping training practical on a
single GPU. Future work will extend the protocol to larger datasets and stronger
open-vocabulary backbones.

\bibliographystyle{plain}
\bibliography{references}
\end{document}
